\documentclass[
    11pt,
    a4paper,
    parskip=half,
    bibliography=totoc,
    numbers=noenddot
]{scrartcl}

%================================================
% Typography
%================================================

\usepackage[utf8]{inputenc}
\usepackage[T1]{fontenc}

% UniGe thesis-family typography:
% Latin Modern keeps the classic Computer Modern appearance used by
% standard thesis classes; FiraMono matches the monospaced font explicitly
% loaded by the UniGe masterthesis template.
\usepackage{lmodern}
\usepackage{FiraMono}
\usepackage[protrusion=true,expansion=true]{microtype}

%================================================
% Mathematics, figures, tables
%================================================

\usepackage{graphicx}
\usepackage{amsmath}
\usepackage{amssymb}
\usepackage{booktabs}
\usepackage{array}
\usepackage{colortbl}
\usepackage{enumitem}
\usepackage[section]{placeins}
\usepackage{float}
\usepackage{caption}

%================================================
% Colour system
%================================================

\usepackage{xcolor}
\definecolor{Navy}{HTML}{18324A}
\definecolor{Accent}{HTML}{2B6F8A}
\definecolor{AccentSoft}{HTML}{EAF3F6}
\definecolor{TableHeader}{HTML}{EEF4F7}
\definecolor{TextMain}{HTML}{20262C}
\definecolor{TextMuted}{HTML}{68737D}
\definecolor{RuleGray}{HTML}{D5DDE2}

\color{TextMain}

%================================================
% Page geometry
%================================================

\usepackage{geometry}
\geometry{
    top=2.35cm,
    bottom=2.45cm,
    inner=2.65cm,
    outer=2.35cm,
    headheight=15pt,
    headsep=0.75cm
}

\raggedbottom

%================================================
% Running headers and page numbers
%================================================

\usepackage{scrlayer-scrpage}
\clearpairofpagestyles
\automark{section}

\ihead{\sffamily\footnotesize\color{TextMuted} AI4EO Research Landscape}
\ohead{\sffamily\footnotesize\color{TextMuted} \headmark}
\cfoot{\sffamily\small\color{TextMuted} \pagemark}

\setkomafont{pageheadfoot}{\normalfont}
\KOMAoptions{headsepline=0.35pt}
\addtokomafont{headsepline}{\color{RuleGray}}

%================================================
% Section hierarchy
%================================================

\setcounter{secnumdepth}{3}
\setcounter{tocdepth}{2}

\setkomafont{section}{\sffamily\bfseries\Large\color{Navy}}
\setkomafont{subsection}{\sffamily\bfseries\large\color{Navy}}
\setkomafont{subsubsection}{\sffamily\bfseries\normalsize\color{Accent}}
\setkomafont{paragraph}{\sffamily\bfseries\normalsize\color{Navy}}
\setkomafont{disposition}{\normalcolor}

\renewcommand*\sectionformat{%
    \makebox[2.2em][l]{\color{Accent}\thesection\autodot}%
}
\renewcommand*\subsectionformat{%
    \makebox[2.8em][l]{\color{Accent}\thesubsection\autodot}%
}
\renewcommand*\subsubsectionformat{%
    \makebox[3.5em][l]{\color{TextMuted}\thesubsubsection\autodot}%
}

\RedeclareSectionCommand[
    beforeskip=2.4ex plus 0.8ex minus 0.4ex,
    afterskip=1.15ex plus 0.25ex
]{section}

\RedeclareSectionCommand[
    beforeskip=1.8ex plus 0.6ex minus 0.3ex,
    afterskip=0.7ex plus 0.2ex
]{subsection}

\RedeclareSectionCommand[
    beforeskip=1.4ex plus 0.5ex minus 0.2ex,
    afterskip=0.45ex plus 0.15ex
]{subsubsection}

%================================================
% Lists
%================================================

\setlist[itemize]{
    leftmargin=1.7em,
    itemsep=0.18em,
    topsep=0.35em,
    parsep=0pt
}

%================================================
% Figures and captions
%================================================

\captionsetup{
    font=small,
    labelfont={bf,sf,color=Accent},
    textfont={color=TextMain},
    labelsep=period,
    skip=7pt
}

%================================================
% Tables
%================================================

\renewcommand{\arraystretch}{1.16}
\setlength{\tabcolsep}{5.5pt}

%================================================
% Hyperlinks
%================================================

\usepackage{xurl}
\usepackage{hyperref}

\hypersetup{
    colorlinks=true,
    linkcolor=Navy,
    citecolor=Accent,
    urlcolor=Accent,
    pdfauthor={Stefano Infusini},
    pdftitle={Artificial Intelligence for Earth Observation},
    pdfsubject={Research Landscape, Evaluation Challenges, and Open Problems}
}

%================================================
% Professional title-page styling
%================================================

\makeatletter
\renewcommand{\maketitle}{%
    \begin{titlepage}
        \thispagestyle{empty}
        \vspace*{1.45cm}

        {\color{Accent}\rule{2.4cm}{2.2pt}}\par
        \vspace{1.35cm}

        {\sffamily\bfseries\fontsize{29}{34}\selectfont\color{Navy}
            \@title\par}

        \vfill

        {\color{RuleGray}\rule{\textwidth}{0.7pt}}\par
        \vspace{0.6cm}

        {\sffamily\large\color{Navy}\@author\par}
        \vspace{0.18cm}
        {\sffamily\small\color{TextMuted}\@date\par}

        \vspace{0.75cm}
    \end{titlepage}
}
\makeatother

%================================================
% Title
%================================================

\title{
    \textbf{Artificial Intelligence for Earth Observation}\\
    \vspace{0.3cm}
    \large Research Landscape, Evaluation Challenges, and Open Problems
}

\author{Stefano Infusini}
\date{August 2026}

%================================================
% Document
%================================================

\begin{document}

\maketitle
\tableofcontents
\newpage

%================================================
\section{Introduction}

Earth Observation (EO) provides systematic information about the Earth's
surface, oceans, and atmosphere through remote-sensing technologies such
as optical imaging, radar, LiDAR, and related geospatial products.
Modern satellite missions generate large and heterogeneous archives
characterized by different spatial, spectral, temporal, and radiometric
resolutions.

Artificial Intelligence (AI) is now applied to a broad range of EO
problems, including scene classification, land-cover mapping, object
detection, change detection, geophysical parameter estimation,
spatio-temporal analysis, and environmental forecasting. More recently,
research has expanded toward self-supervised learning, multimodal
representation learning, large-scale pretraining, and Earth Observation
foundation models.

The objective of this report is not to provide a complete textbook
treatment of remote sensing or to preselect a thesis topic. Instead, it
aims to provide a compact technical basis for discussing possible MSc
thesis directions. The report therefore focuses on three questions:
which characteristics of EO data are most relevant to AI, which
downstream tasks define the main problem families, and which research
problems are currently receiving significant attention in the AI4EO
literature.

Detailed mathematical formulations and additional remote-sensing
background are collected in the appendices so that the main text remains
focused on the research landscape and on the issues that are most
relevant when evaluating possible thesis directions.

%================================================
\section{Earth Observation for AI}

\subsection{EO Data Modalities}

Earth Observation data differ substantially from conventional RGB
computer-vision datasets. Depending on the sensor, the input may contain
a few visible channels, many spectral bands, radar backscatter,
three-dimensional measurements, elevation data, or combinations of
multiple modalities. These differences directly affect model design,
preprocessing, pretraining, data fusion, and evaluation.

\subsubsection{Optical, Multispectral, and Hyperspectral Imagery}

Optical remote sensing measures reflected solar radiation and includes
conventional RGB, multispectral, and hyperspectral imaging. Different
materials exhibit wavelength-dependent spectral responses, allowing
vegetation, water, soil, snow, rock, and artificial surfaces to be
distinguished through their spectral behaviour
\cite{esa_optical_imagery,esa_reflectance_curves}.

\begin{figure}[H]
    \centering
    \IfFileExists{images/reflectance_plot.jpg}{
        \includegraphics[
            width=0.95\linewidth,
            height=0.38\textheight,
            keepaspectratio
        ]{images/reflectance_plot.jpg}
    }{
        \fbox{
            \parbox[c][4.8cm][c]{0.88\linewidth}{
                \centering
                Add \texttt{images/reflectance\_plot.jpg}
            }
        }
    }
    \caption{
        Representative spectral reflectance curves of different surface
        materials and atmospheric transmission windows. Source: ESA
        \cite{esa_reflectance_curves}.
    }
    \label{fig:reflectance-curves-earth-features}
\end{figure}

Multispectral sensors acquire a limited number of broad spectral bands.
Sentinel-2, for example, acquires 13 bands spanning visible, red-edge,
near-infrared, and short-wave infrared regions, with native spatial
resolutions of 10, 20, and 60 m \cite{esa_sentinel2_msi}. This is
important for AI because a model may receive more than three channels,
and because bands may have to be aligned or resampled before being used
jointly.

Hyperspectral sensors acquire tens or hundreds of narrow and often
contiguous bands. Their dense spectral sampling can reveal absorption
features that are not visible in RGB or conventional multispectral
imagery, but also produces high-dimensional inputs, strong correlation
between neighbouring bands, larger storage requirements, and a greater
risk of overfitting when labels are scarce
\cite{jpl_imaging_spectroscopy,bioucasdias2013hyperspectral}.

For AI4EO, the key distinction is therefore not simply the number of
channels. Optical EO data may involve heterogeneous band resolutions,
spectral redundancy, atmospheric effects, mixed pixels, missing
observations due to clouds, and sensor-dependent spectral responses.

\subsubsection{Synthetic Aperture Radar}

Synthetic Aperture Radar (SAR) is an active microwave sensing technology.
The sensor transmits electromagnetic pulses toward the Earth's surface
and measures the energy scattered back toward the instrument
\cite{moreira2013sar}. Because SAR provides its own illumination and
operates at microwave wavelengths, it can acquire observations during
day or night and can generally observe through cloud cover
\cite{esa_sentinel1_instrument}.

SAR imagery differs fundamentally from optical imagery. The measured
response depends on wavelength, incidence angle, polarization, surface
roughness, dielectric properties, moisture, geometry, and scattering
mechanisms. Smooth water often appears dark because most of the signal
is reflected away from the sensor, while rough surfaces, vegetation, and
urban structures may produce stronger backscatter.

SAR is also a coherent imaging system and is affected by speckle.
Furthermore, its side-looking geometry can produce foreshortening,
layover, and radar shadow. These properties make SAR an important
example of why EO cannot always be treated as ordinary natural-image
computer vision.

Sentinel-1 is a representative spaceborne SAR mission. Its data are used
for flood mapping, soil-moisture analysis, forest monitoring, sea-ice
observation, maritime surveillance, and surface-deformation analysis.
Interferometric SAR (InSAR) additionally exploits phase differences
between multiple SAR acquisitions.

\subsubsection{LiDAR and Elevation Data}

LiDAR is an active optical technology that emits laser pulses and
measures their return time to estimate three-dimensional structure
\cite{baltsavias1999airborne}. LiDAR measurements are commonly
represented as point clouds and can be used to characterize terrain,
vegetation structure, buildings, and other three-dimensional objects.

Elevation information may also be represented through Digital Elevation
Models (DEMs), Digital Surface Models (DSMs), and Digital Terrain Models
(DTMs). These are derived geospatial products rather than sensors
themselves \cite{guth2021dem,usgs_dem_definition}. They can be used as
additional model inputs, for example to provide terrain context for
classification, segmentation, hydrology, landslide assessment, or
environmental modelling.

From an AI perspective, LiDAR and elevation products introduce additional
forms of heterogeneity: irregular point sets, rasterized height products,
different horizontal and vertical accuracies, and the need to combine
three-dimensional or topographic information with optical or radar data.

\subsection{Spatial, Spectral, Temporal, and Radiometric Resolution}

EO sensors are commonly characterized through four forms of resolution.
Spatial resolution determines the ground area represented by a pixel.
Spectral resolution describes how finely the electromagnetic spectrum is
sampled. Temporal resolution describes how frequently the same area can
be observed, while radiometric resolution describes sensitivity to
differences in measured energy.

These dimensions are especially important in AI4EO because datasets from
different sensors rarely share the same sampling characteristics.
Combining Sentinel-1, Sentinel-2, Landsat, LiDAR, DEMs, or commercial
very-high-resolution imagery may therefore require co-registration,
resampling, temporal matching, cloud filtering, and careful treatment of
the information lost or interpolated during preprocessing.

The main consequence is that ``more data'' does not automatically imply
more usable information. Spatial, spectral, and temporal resolutions are
often traded against each other, and the model must operate on
measurements generated by different physical processes and acquisition
geometries.

\subsection{EO Data Availability and Access}

Data accessibility is an important practical constraint for AI4EO
research. A scientifically interesting problem may still be unsuitable
for a thesis if the relevant data are commercial, geographically
restricted, difficult to reproduce, or too expensive to process at the
required scale.

\begin{table}[H]
    \centering
    \small
    \begin{tabular}{
        p{0.20\textwidth}
        p{0.27\textwidth}
        p{0.25\textwidth}
        p{0.17\textwidth}}
        \toprule
        \textbf{Source} &
        \textbf{Representative Data} &
        \textbf{Programmatic Access} &
        \textbf{Access Model} \\
        \midrule

        Copernicus Data Space Ecosystem &
        Sentinel missions and other Copernicus collections &
        STAC, OData, Sentinel Hub and related APIs &
        Open \\

        USGS / EROS &
        Landsat and other EROS archive products &
        Machine-to-Machine (M2M) API &
        Open \\

        NASA Earthdata &
        MODIS, ICESat-2, SMAP and other Earth-science products &
        Earthdata services and \texttt{earthaccess} &
        Open \\

        Planet &
        PlanetScope and other high-resolution imagery &
        Data API and Processing API &
        Commercial / account-dependent \\

        \bottomrule
    \end{tabular}
    \caption{Representative data-access ecosystems relevant to AI4EO research.}
    \label{tab:eo-data-access}
\end{table}

Open archives such as the Copernicus Data Space Ecosystem, USGS
EarthExplorer/EROS, and NASA Earthdata make it possible to construct
reproducible acquisition pipelines
\cite{cdse_documentation,cdse_apis,usgs_m2m,nasa_earthdata_search,nasa_earthaccess}.
Commercial providers can offer higher spatial or temporal resolution,
but licensing and cost can limit reproducibility
\cite{planet_data_api,planet_processing_api}.

Python clients previously developed for programmatic access to CDSE and
the USGS M2M API are publicly available at:

\begin{itemize}
    \item \url{https://github.com/stefanoIn/CDSE-api}
    \item \url{https://github.com/stefanoIn/USGS-m2m-client}
\end{itemize}

%================================================
\section{AI Tasks in Earth Observation}

The task categories below describe the type of mapping learned from EO
data rather than a specific application domain. A single application can
combine multiple task families. Wildfire assessment, for example, may
involve change detection, segmentation, regression, and time-series
analysis. Compact mathematical formulations are included directly within the
corresponding task descriptions.

\subsection{Classification}

Image or scene classification assigns a semantic label to an entire EO
image, image patch, parcel, or region. Typical examples include aerial
scene classification, crop-type classification, and assigning a dominant
land-use category to a spatial unit. Dense pixel-level land-cover mapping
is instead a segmentation problem.

The AID benchmark is a representative aerial scene-classification
dataset containing categories such as farmland, forest, industrial
areas, residential areas, and airports \cite{xia2017aid}.

For an input image

\[
\mathbf{X}\in\mathbb{R}^{C\times H\times W},
\]

classification can be written as

\[
f:
\mathbb{R}^{C\times H\times W}
\rightarrow
\{1,\ldots,K\}.
\]

\begin{figure}[H]
    \centering
    \IfFileExists{images/aid_classification.png}{
        \includegraphics[
            width=0.86\linewidth,
            height=0.34\textheight,
            keepaspectratio
        ]{images/aid_classification.png}
    }{
        \fbox{
            \parbox[c][4.5cm][c]{0.84\linewidth}{
                \centering
                Add \texttt{images/aid\_classification.png}
            }
        }
    }
    \caption{
        Example images from the AID dataset, illustrating aerial
        scene classification. Reproduced from \cite{peng2022mopnet}.
    }
    \label{fig:aid-classification}
\end{figure}

\subsection{Semantic and Instance Segmentation}

Semantic segmentation assigns a semantic class to every pixel and is
used for tasks such as land-cover mapping, flood delineation,
burned-area mapping, road extraction, and vegetation mapping.

Instance segmentation additionally separates distinct objects belonging
to the same class. This is useful when individual buildings, vehicles,
ships, or tree crowns must be identified rather than only assigning a
shared semantic category to their pixels.

For semantic segmentation, the mapping can be written as

\[
f:
\mathbb{R}^{C\times H\times W}
\rightarrow
\{1,\ldots,K\}^{H\times W}.
\]

For instance segmentation, a simplified mask representation is

\[
f:
\mathbb{R}^{C\times H\times W}
\rightarrow
\{0,1\}^{D\times H\times W},
\]

where \(D\) is the image-dependent number of detected instances.

\begin{figure}[H]
    \centering
    \IfFileExists{images/semantic_segmentation.png}{
        \includegraphics[
            width=0.90\linewidth,
            height=0.32\textheight,
            keepaspectratio
        ]{images/semantic_segmentation.png}
    }{
        \fbox{
            \parbox[c][4.2cm][c]{0.84\linewidth}{
                \centering
                Add \texttt{images/semantic\_segmentation.png}
            }
        }
    }
    \caption{
        Example of semantic segmentation applied to aerial imagery,
        showing the input image and the corresponding pixelwise semantic
        mask. Adapted from \cite{davies2022semantic}.
    }
    \label{fig:semantic-segmentation}
\end{figure}

\subsection{Object Detection}

Object detection identifies and localizes individual objects using
bounding regions together with class labels and confidence scores. In
very-high-resolution EO imagery it is commonly applied to buildings,
vehicles, ships, aircraft, and infrastructure.

The localization component can be represented as

\[
f:
\mathbb{R}^{C\times H\times W}
\rightarrow
\mathbb{R}^{D\times 4},
\]

with additional class labels and confidence scores in a complete detector.

\begin{figure}[H]
    \centering
    \IfFileExists{images/object_detection.jpg}{
        \includegraphics[
            width=0.86\linewidth,
            height=0.36\textheight,
            keepaspectratio
        ]{images/object_detection.jpg}
    }{
        \fbox{
            \parbox[c][4.5cm][c]{0.84\linewidth}{
                \centering
                Add \texttt{images/object\_detection.jpg}
            }
        }
    }
    \caption{
        Example of object detection in high-resolution remote-sensing
        imagery, where individual objects are localized using bounding
        boxes. Reproduced from \cite{ballinger2023rs4osint}.
    }
    \label{fig:object-detection}
\end{figure}

\subsection{Change Detection}

Change detection identifies meaningful differences between observations
of the same area acquired at different times. Applications include urban
expansion, deforestation, flood mapping, wildfire assessment,
infrastructure monitoring, and disaster response.

A central EO-specific difficulty is that image differences do not always
correspond to real changes on the ground. Illumination, atmospheric
conditions, seasonality, viewing geometry, sensor characteristics,
misregistration, and cloud cover can all create apparent differences.
Change-detection methods must therefore distinguish semantic change from
acquisition-related variability.

For two co-registered acquisitions,

\[
\mathbf{X}_{t_1},\mathbf{X}_{t_2}
\in
\mathbb{R}^{C\times H\times W}.
\]

Binary change detection may be written as

\[
f:
\mathbb{R}^{2\times C\times H\times W}
\rightarrow
\{0,1\}^{H\times W},
\]

while semantic change detection may use

\[
f:
\mathbb{R}^{2\times C\times H\times W}
\rightarrow
\{1,\ldots,K\}^{H\times W}.
\]

\subsection{Regression and Geophysical Parameter Estimation}

Regression tasks estimate continuous-valued quantities from EO data.
Predictions may correspond to an entire region or to a spatially
distributed continuous field. Typical applications include biomass
estimation, land-surface temperature estimation, soil-moisture retrieval,
air-quality variables, and other biophysical or geophysical parameters.

Unlike segmentation, where each spatial location is assigned a discrete
class, pixelwise regression assigns a continuous value to each spatial
location.

Image-level regression can be written as

\[
f:
\mathbb{R}^{C\times H\times W}
\rightarrow
\mathbb{R},
\]

while pixelwise regression is

\[
f:
\mathbb{R}^{C\times H\times W}
\rightarrow
\mathbb{R}^{H\times W}.
\]

\subsection{Time-Series Analysis and Forecasting}

Satellite image time series consist of repeated observations of the same
geographical area over time. The temporal dimension can support
classification, change or event detection, regression, trend analysis,
and forecasting.

Applications include vegetation and crop monitoring, air-quality
monitoring, land-cover dynamics, disaster monitoring, and weather or
climate forecasting. Important EO-specific challenges include irregular
acquisition times, missing observations, cloud contamination, seasonal
variability, long temporal dependencies, and differences in revisit time
between sensors.

A satellite image time series may be represented as

\[
\mathbf{X}
\in
\mathbb{R}^{T\times C\times H\times W}.
\]

A spatial forecasting problem can then be written as

\[
f:
\mathbb{R}^{T\times C\times H\times W}
\rightarrow
\mathbb{R}^{C'\times H\times W}.
\]

\subsection{Image Reconstruction and Enhancement}

Image reconstruction and enhancement aim to recover or improve an EO
image or continuous EO product rather than directly predicting a
semantic label. Typical applications include super-resolution,
statistical or learned downscaling, pansharpening, denoising, cloud
removal, image fusion, and reconstruction of missing observations.

These problems are especially relevant in EO because different sensors
often trade spatial, spectral, and temporal resolution against one
another.

A general reconstruction mapping is

\[
f:
\mathbb{R}^{C_{\mathrm{in}}\times H_{\mathrm{in}}\times W_{\mathrm{in}}}
\rightarrow
\mathbb{R}^{C_{\mathrm{out}}\times H_{\mathrm{out}}\times W_{\mathrm{out}}}.
\]

For super-resolution,

\[
f:
\mathbb{R}^{C\times H\times W}
\rightarrow
\mathbb{R}^{C\times sH\times sW}.
\]

\subsection{Vision-Language Tasks}

Vision-language tasks connect EO imagery with natural language.
Representative tasks include image captioning, visual question answering,
cross-modal retrieval, text-image matching, and grounding.

In EO, these tasks are becoming increasingly relevant because they can
connect large image archives with semantic descriptions and may support
natural-language interaction with geospatial data. They also form part
of the broader transition toward multimodal and agentic geospatial
systems discussed later in the report.

For a vocabulary of size \(V\) and a caption of length \(L\), a captioning
output can be represented as

\[
f:
\mathbb{R}^{C\times H\times W}
\rightarrow
\{1,\ldots,V\}^{L}.
\]

%================================================
\section{Current AI4EO Research Landscape}

The methodological evolution of AI4EO broadly follows the development of
machine learning itself: hand-crafted features were progressively
replaced by learned representations, supervised deep learning was
followed by self-supervised and large-scale pretraining, and
Transformer-based models increasingly enabled large reusable
representations. The important question for this report is therefore not
the chronology itself, but how these developments interact with the
specific properties of EO data.

The current landscape is shaped by recurring EO-specific problems:
large unlabeled archives, expensive and geographically limited labels,
heterogeneous sensors, different native resolutions, irregular temporal
sampling, distribution shift across geography and season, and the need
to evaluate models under realistic operational conditions.

\subsection{Foundation Models and Large-Scale Pretraining}

% To be populated from the literature review.
% Focus on the move from task-specific models to reusable pretrained
% representations, self-supervised pretraining, masked modelling,
% adaptation strategies, sensor-specific vs sensor-agnostic models,
% and whether large-scale pretraining consistently improves downstream
% performance.

\subsection{Multimodal and Multisensor Learning}

% To be populated from the literature review.
% Optical + SAR + hyperspectral + LiDAR/DEM + meteorological variables +
% in-situ observations + text.
% Discuss co-registration, asynchronous acquisitions, resolution mismatch,
% missing modalities, and fusion strategies.

\subsection{Spatio-Temporal Learning}

% To be populated from the literature review.
% Satellite image time series, bi-temporal and multi-temporal learning,
% forecasting, event evolution, irregular sampling, cloud-related gaps,
% long temporal context, and seasonality.

\subsection{Generalization, Domain Shift, and Limited Supervision}

% To be populated from the literature review.
% Geographic, seasonal, sensor, and resolution shift.
% Domain adaptation, semi-supervision, weak supervision, few-shot learning,
% active learning, OOD detection, and evaluation outside the training domain.

\subsection{Trustworthy and Uncertainty-Aware AI}

% To be populated from the literature review.
% Uncertainty quantification, calibration, explainability, robustness,
% failure detection, out-of-distribution behaviour, and scientific reliability.

\subsection{Physics-Aware and Hybrid AI}

% To be populated from the literature review.
% Physical constraints, differentiable physical models, hybrid
% physical/data-driven approaches, radiative-transfer knowledge,
% conservation laws, and scientific priors.

\subsection{Vision-Language and Agentic Geospatial AI}

% To be populated from the literature review.
% Image-text representation learning, captioning, retrieval, grounding,
% VLMs, domain-specific LLMs, RAG, tool use, and agentic workflows.
% Distinguish vision-language representation learning from LLM-based assistants.

\subsection{Efficient and Operational AI}

% To be populated from the literature review.
% Parameter-efficient adaptation, distillation, compression, cloud/HPC,
% scalable inference, edge/onboard AI, data-transfer constraints,
% and operational deployment.

%================================================
\section{Evaluation, Benchmarks, and Open Problems}

Evaluation is a central issue in AI4EO because performance measured on a
single benchmark may not reflect geographic transfer, sensor changes,
limited-label settings, or operational cost. This section should
therefore examine not only which models perform best, but also how
comparisons are constructed and which dimensions of performance remain
poorly measured.

\subsection{Benchmarking AI4EO Models}

% Discuss EO task suites and benchmarking initiatives such as PANGAEA,
% GEO-Bench and related work.
% Compare what they measure: tasks, modalities, resolutions, data regimes,
% adaptation protocols, baselines, and computational requirements.
% The purpose is to assess what current benchmarks reveal and what they miss,
% not to assume that a specific benchmark or model family will become the thesis.

\subsection{Geographic, Temporal, and Sensor Generalization}

% Geographic train/test separation, seasonal and temporal transfer,
% sensor transfer, resolution changes, acquisition-condition shifts,
% and the distinction between in-distribution benchmark performance
% and deployment robustness.

\subsection{Label Efficiency and Data Availability}

% Limited labels, annotation cost, weak/noisy labels, few-shot settings,
% self-supervised pretraining, active learning, and whether models remain useful
% when labelled data are scarce.

\subsection{Multimodal and Multi-Resolution Integration}

% Co-registration, asynchronous observations, missing modalities,
% heterogeneous native resolutions, different measurement physics,
% and the effect of resampling on fair comparisons.

\subsection{Computational Cost and Scalability}

% Training/fine-tuning cost, GPU memory, inference cost, storage,
% data transfer, model size, cloud/HPC requirements, and whether
% performance gains justify additional computational cost.

\subsection{Reproducibility and Fair Evaluation}

% Dataset leakage, geographic leakage, temporal leakage, sensor leakage,
% hyperparameter budgets, preprocessing differences, baseline strength,
% code/data availability, and reproducible protocols.

%================================================
\section{Discussion}

The discussion should synthesize the literature rather than introduce a
preselected thesis topic. In particular, it should identify which open
problems recur across multiple branches of AI4EO, which claimed
advantages of modern approaches are well supported, and which remain
dependent on benchmark design, data availability, or computational
resources.

The final version of this section can also compare the practical
feasibility of different research questions in terms of available data,
compute, implementation effort, evaluation clarity, and scientific
relevance. Possible thesis directions should emerge from this synthesis
only after the research landscape and benchmark evidence have been
reviewed.

%================================================
\section{Conclusion}

The conclusion should summarize the principal characteristics that make
AI4EO distinct from conventional computer vision, the strongest current
research directions, and the main evaluation and feasibility challenges.
It should provide a basis for discussion with the supervisors without
preselecting a specific thesis topic.

%================================================
% Appendices
%================================================

\appendix

\section{Technical Notes on Earth Observation Modalities}
\label{app:eo-technical}

This appendix collects mathematical details that are useful for technical
reference but are not necessary for the main research-landscape
discussion.

\subsection{Optical and Multispectral Representations}

A multispectral image aligned to a common spatial grid can be represented
as

\[
\mathbf{X}\in\mathbb{R}^{H\times W\times B},
\]

where \(H\) and \(W\) are spatial dimensions and \(B\) is the number of
bands. A pixel \(p\) is represented by

\[
\mathbf{x}_p =
\left[
x_{p,1},
x_{p,2},
\ldots,
x_{p,B}
\right].
\]

A spectral band integrates incoming radiance over a finite wavelength
range. A simplified measurement model is

\[
x_{p,b}
=
\frac{
    \int L_p(\lambda)S_b(\lambda)\,d\lambda
}{
    \int S_b(\lambda)\,d\lambda
},
\]

where \(L_p(\lambda)\) is the spectral radiance reaching the sensor and
\(S_b(\lambda)\) is the spectral response function of band \(b\).

For Sentinel-2, a pixel may be represented by the 13-band vector

\[
\mathbf{x}_p =
\left[
B_1,B_2,B_3,B_4,B_5,B_6,B_7,B_8,B_{8A},B_9,B_{10},B_{11},B_{12}
\right].
\]

A true-colour composite is commonly formed as

\[
\mathrm{RGB}_{\mathrm{true}}=[B_4,B_3,B_2],
\]

while a common false-colour composite is

\[
\mathrm{RGB}_{\mathrm{false}}=[B_8,B_4,B_3].
\]

The Normalized Difference Vegetation Index is

\[
\mathrm{NDVI}
=
\frac{\rho_{\mathrm{NIR}}-\rho_{\mathrm{Red}}}
     {\rho_{\mathrm{NIR}}+\rho_{\mathrm{Red}}}
\]

\cite{rouse1974monitoring}.

\subsection{Hyperspectral Representations}

Hyperspectral imagery is also represented as a data cube,

\[
\mathbf{X}\in\mathbb{R}^{H\times W\times B},
\]

but \(B\) may contain tens or hundreds of narrow bands. The bandwidth of
band \(b\) is

\[
\Delta\lambda_b =
\lambda_{b,\max}-\lambda_{b,\min},
\]

and the corresponding resolving power can be written as

\[
R_b=\frac{\lambda_b}{\Delta\lambda_b}.
\]

In a simplified linear spectral-mixture model,

\[
\mathbf{x}_p
=
\sum_{k=1}^{K}a_{p,k}\mathbf{s}_k
+
\boldsymbol{\epsilon}_p,
\]

subject to

\[
a_{p,k}\geq0,
\qquad
\sum_{k=1}^{K}a_{p,k}=1.
\]

Here, \(\mathbf{s}_k\) represents an endmember spectrum and \(a_{p,k}\)
its fractional abundance within pixel \(p\).

\subsection{SAR Measurements}

The radar range can be estimated from the round-trip travel time:

\[
R=\frac{c\Delta t}{2}.
\]

A focused Single Look Complex SAR pixel can be written as

\[
S(p)=I(p)+jQ(p)=A(p)e^{j\phi(p)},
\]

with intensity

\[
I_{\mathrm{SAR}}(p)=|S(p)|^2.
\]

After radiometric calibration, backscatter is commonly expressed as

\[
\sigma^0_{\mathrm{dB}}
=
10\log_{10}(\sigma^0).
\]

A simplified representation of factors affecting the measured
backscatter is

\[
\sigma^0 =
f(\lambda,\theta,P,\varepsilon,r,g),
\]

where the variables denote wavelength, incidence angle, polarization,
dielectric properties, surface roughness, and observation geometry.

The principal polarization channels are

\[
HH,\qquad HV,\qquad VH,\qquad VV.
\]

A dual-polarization Sentinel-1 observation may therefore be represented as

\[
\mathbf{x}_p =
\left[
\sigma^0_{VV}(p),
\sigma^0_{VH}(p)
\right].
\]

\subsection{LiDAR and Elevation}

The LiDAR range equation is

\[
R=\frac{c\Delta t}{2}.
\]

A LiDAR return may be represented as

\[
\mathbf{p}_i=
[x_i,y_i,z_i,I_i,r_i],
\]

where the first three components are spatial coordinates and the
remaining components may encode return intensity and return number.

For elevation products,

\[
D(i,j)=z_{i,j},
\]

and a normalized Digital Surface Model can be computed as

\[
\mathrm{nDSM}(i,j)=\mathrm{DSM}(i,j)-\mathrm{DTM}(i,j).
\]


%================================================
% Bibliography
%================================================

\begin{thebibliography}{99}

\bibitem{esa_optical_imagery}
European Space Agency,
\textit{Optical Earth Observation},
ESA Earth Observation.
Available online:
\url{https://www.esa.int/Applications/Observing_the_Earth}

\bibitem{esa_reflectance_curves}
European Space Agency,
\textit{Reflectance Curves of Snow, Vegetation, Water and Rock},
ESA Eduspace.
Available online:
\url{https://www.esa.int/SPECIALS/Eduspace_EN/}

\bibitem{esa_sentinel2_msi}
European Space Agency,
\textit{Sentinel-2 Multispectral Instrument Overview},
ESA Sentinel Application Platform.
Available online:
\url{https://step.esa.int/main/wp-content/help/versions/11.0.0/snap-toolboxes/eu.esa.opt.opttbx.s2msi.reader/Sentinel2Overview.html}

\bibitem{jpl_imaging_spectroscopy}
NASA Jet Propulsion Laboratory,
\textit{AVIRIS Imaging Spectroscopy}.
Available online:
\url{https://aviris.jpl.nasa.gov/aviris/imaging_spectroscopy.html}

\bibitem{bioucasdias2013hyperspectral}
J.~M. Bioucas-Dias, A.~Plaza, G.~Camps-Valls, P.~Scheunders,
N.~M. Nasrabadi, and J.~Chanussot,
``Hyperspectral remote sensing data analysis and future challenges,''
\textit{IEEE Geoscience and Remote Sensing Magazine},
vol.~1, no.~2, pp.~6--36, 2013.

\bibitem{rouse1974monitoring}
J.~W. Rouse, R.~H. Haas, J.~A. Schell, and D.~W. Deering,
``Monitoring vegetation systems in the Great Plains with ERTS,''
in \textit{Proceedings of the Third Earth Resources Technology
Satellite-1 Symposium},
vol.~1, pp.~309--317, 1974.

\bibitem{moreira2013sar}
A.~Moreira, P.~Prats-Iraola, M.~Younis, G.~Krieger, I.~Hajnsek,
and K.~P. Papathanassiou,
``A tutorial on synthetic aperture radar,''
\textit{IEEE Geoscience and Remote Sensing Magazine},
vol.~1, no.~1, pp.~6--43, 2013.

\bibitem{esa_sentinel1_instrument}
European Space Agency,
\textit{Sentinel-1 Instrument},
ESA Earth Observation.
Available online:
\url{https://www.esa.int/Applications/Observing_the_Earth/Copernicus/Sentinel-1/Instrument}

\bibitem{baltsavias1999airborne}
E.~P. Baltsavias,
``Airborne laser scanning: Basic relations and formulas,''
\textit{ISPRS Journal of Photogrammetry and Remote Sensing},
vol.~54, no.~2--3, pp.~199--214, 1999.

\bibitem{usgs_lidar_data}
United States Geological Survey,
\textit{What Is Lidar Data and Where Can I Download It?}
Available online:
\url{https://www.usgs.gov/faqs/what-lidar-data-and-where-can-i-download-it}

\bibitem{guth2021dem}
P.~L. Guth, A.~Van Niekerk, C.~H. Grohmann, J.-P. Muller,
L.~Hawker, I.~V. Florinsky, D.~B. Gesch, H.~I. Reuter,
V.~Herrera-Cruz, S.~Riazanoff, C.~Lopez-Vazquez,
C.~C. Carabajal, C.~Albinet, and P.~Strobl,
``Digital elevation models: Terminology and definitions,''
\textit{Remote Sensing},
vol.~13, no.~18, Art.~no.~3581, 2021.

\bibitem{usgs_dem_definition}
United States Geological Survey,
\textit{What Is a Digital Elevation Model (DEM)?}
Available online:
\url{https://www.usgs.gov/faqs/what-a-digital-elevation-model-dem}

\bibitem{esa_earth_observation}
European Space Agency,
\textit{Newcomers Earth Observation Guide}.
Available online:
\url{https://business.esa.int/newcomers-earth-observation-guide}

\bibitem{cdse_documentation}
Copernicus Data Space Ecosystem,
\textit{Copernicus Data Space Ecosystem Documentation}.
Available online:
\url{https://documentation.dataspace.copernicus.eu/}

\bibitem{cdse_apis}
Copernicus Data Space Ecosystem,
\textit{APIs}.
Available online:
\url{https://documentation.dataspace.copernicus.eu/APIs.html}

\bibitem{cdse_browser}
Copernicus Data Space Ecosystem,
\textit{About the Copernicus Browser}.
Available online:
\url{https://documentation.dataspace.copernicus.eu/Applications/Browser.html}

\bibitem{usgs_m2m}
United States Geological Survey,
\textit{Machine-to-Machine (M2M) API}.
Available online:
\url{https://m2m.cr.usgs.gov/}

\bibitem{nasa_earthdata_search}
National Aeronautics and Space Administration,
\textit{Earthdata Search}.
Available online:
\url{https://www.earthdata.nasa.gov/data/tools/earthdata-search}

\bibitem{nasa_earthaccess}
National Aeronautics and Space Administration,
\textit{earthaccess}.
Available online:
\url{https://www.earthdata.nasa.gov/data/tools/earthaccess}

\bibitem{planet_data_api}
Planet Labs,
\textit{Data API Overview}.
Available online:
\url{https://docs.planet.com/develop/apis/data/}

\bibitem{planet_processing_api}
Planet Labs,
\textit{Processing API}.
Available online:
\url{https://docs.planet.com/develop/apis/processing/}

\bibitem{peng2022mopnet}
F.~Peng, W.~Lu, W.~Tan, K.~Qi, X.~Zhang, and Q.~Zhu,
``Multi-Output Network Combining GNN and CNN for Remote Sensing Scene Classification,''
\textit{Remote Sensing},
vol.~14, no.~6, Art.~no.~1478, 2022.

\bibitem{xia2017aid}
G.-S. Xia, J.~Hu, F.~Hu, B.~Shi, X.~Bai, Y.~Zhong, L.~Zhang,
and X.~Lu,
``AID: A Benchmark Data Set for Performance Evaluation of Aerial
Scene Classification,''
\textit{IEEE Transactions on Geoscience and Remote Sensing},
vol.~55, no.~7, pp.~3965--3981, 2017.

\bibitem{ni2024objectcounts}
Z.~Ni, Z.~Zong, and P.~Ren,
``Incorporating object counts into remote sensing image captioning,''
\textit{International Journal of Digital Earth},
vol.~17, no.~1, Art.~no.~2392847, 2024.

\bibitem{davies2022semantic}
A.~J. Davies,
\textit{Semantic Segmentation of Aerial Imagery using U-Net in Python},
Towards Data Science, Mar.~31, 2022.
Available online:
\url{https://towardsdatascience.com/semantic-segmentation-of-aerial-imagery-using-u-net-in-python-552705238514/}

\bibitem{liu2025catnet}
Y.~Liu, H.~Li, C.~Hu, S.~Luo, Y.~Luo, and C.~W. Chen,
``Learning to Aggregate Multi-Scale Context for Instance Segmentation
in Remote Sensing Images,''
\textit{IEEE Transactions on Neural Networks and Learning Systems},
vol.~36, no.~1, pp.~595--609, 2025.

\bibitem{ballinger2023rs4osint}
O.~Ballinger,
\textit{Remote Sensing for OSINT: Object Detection},
2023.
Available online:
\url{https://bellingcat.github.io/RS4OSINT/C5_Object_Detection.html}

\end{thebibliography}

\end{document}
